Maxwell's equations govern the behavior of electromagnetic fields and form the
foundation of classical electromagnetism. These equations describe how electric
and magnetic fields interact and propagate through space. In their differential
form, they are
\begin{align*}
\nabla \cdot \boldsymbol{E} &= \frac{\rho}{\varepsilon_0}\\
\nabla \cdot \boldsymbol{B} &= 0\\
\nabla \times \boldsymbol{E} &= -\frac{\partial \boldsymbol{B}}{\partial t}\\
\nabla \times \boldsymbol{B} &=
\frac{1}{c^2}\cdot\frac{\partial \boldsymbol{E}}{\partial t}
+ \mu_0 \boldsymbol{J},
\end{align*}

This form is known as the microscopic form of Maxwell's equations.
Another form that we will use is called the macroscopic form, this
one is better suited for working with materials parameters, in the
case of vacuum, which is what we care about in this paper we have
\(\varepsilon_0\boldsymbol{E}=\boldsymbol{D}\),
\(\boldsymbol{B}=\mu_0\boldsymbol{H}\) and
\(\mu_0\varepsilon_0=\frac{1}{c^2}\), this leads to:
\begin{align*}
\nabla \cdot \boldsymbol{D} &= \rho\\
\nabla \cdot \boldsymbol{B} &= 0\\
\nabla \times \boldsymbol{E} &= -\frac{\partial \boldsymbol{B}}{\partial t}\\
\nabla \times \boldsymbol{H} &=  \frac{\partial \boldsymbol{D}}{\partial t}
+ \boldsymbol{J}.
\end{align*}
In the case of dielectric material the constants \(\varepsilon_0\)
and \(\mu_0\) are replaced by matrices and the relationships read
\(\varepsilon^{ij}\boldsymbol{E}=\boldsymbol{D}\) and
\(\boldsymbol{B}=\mu^{ij}\boldsymbol{H}\).

We can see that Maxwell's equations involve two key differential
operators: the divergence and the curl. Unlike the scalar Helmholtz
equation, electromagnetic fields are
vector quantities, and their transformation under coordinate
changes involves both the divergence and curl operators.
Understanding how these operators transform is essential for
deriving the material parameters that achieve cloaking.

\subsection{Form invariance of Maxwell's equations}

Something remarkably about Maxwell's equations is that they possess
a special property called form invariance, meaning they maintain
their structure under any coordinate transformation. This property
is so fundamental that according to
\cite{mccallRoadmapTransformationOptics2018} this is one of the
guiding principles Einstein used to develop his theory of general
relativity, it is also the key mathematical principle that enables
transformation optics to work with the full electromagnetic
equations.

\subsection{The transformation}

The next step is to come up with a coordinate transformation that
creates a region of space where we can hide stuff, this region of
space needs to be decoupled from the original space; what does this
mean for the transformation? It basically means that the transformation maps
a point, a line or a small circle into a large circle, which is our
cloaking region, this one will be perceived as very small despite being large.
In this paper the transformation that we will use is described in
terms of cylindrical coordinates as follows:
\[
\begin{cases}
	r= \frac{b-a}{b} r' + a\quad r'\leq b\\
	\theta = \theta'\\
	z=z'
\end{cases},
\]
for \(r'\geq b\) the transformation is the identity.

The primed coordinates describe what is known as the virtual space
which can be thought of as the space that light perceives. The
unprimed coordinates describe the physical space; i.e. the actual
space. Notice that \(0\leq r'\leq b\) gets mapped to \(a\leq r\leq
b\), so nothing is mapped to \(0\leq r \leq a\), this is the space
that will be cloaked, since from the perspective of light it
doesn't exist; or rather, it is perceived as a single point (the
origin).

\subsection{Encoding the transformation into the material parameters}

At this point we could in principle follow the procedure outlined in section
\ref{helmholtzexample}, using the chain rule to re-write the operators in the
equations and from there get the material parameters; however, this will lead
to complicated and long expressions that can be hard to deal with. The
transformation optics approach is to come up with a formula that allows us to
do this in the most general sense, then we can plug in the specific
transformation we want and obtain the parameters from the formula. This
approach requires more sophisticated theory and knowledge about differential
geometry but it is more elegant and simplifies computations a lot.
